\chapter{Introducción}
\label{chap:introduccion}

\lettrine{P}{ython} ocupa una posición consolidada en la computación científica.
Su ecosistema de bibliotecas, la legibilidad del código y la rapidez de
prototipado lo han convertido en el lenguaje de referencia para análisis
numérico y aprendizaje automático~\cite{HighPerfPython2020}. Sin embargo, la
ejecución paralela con múltiples hilos ha estado históricamente limitada por el
\emph{Global Interpreter Lock} (GIL), que impide la ejecución simultánea de
bytecode Python en un mismo proceso. Las versiones \emph{free-threaded} de Python, disponibles
experimentalmente a partir de Python~3.13~\cite{PEP703}, eliminan esa
restricción y abren la posibilidad de explotar paralelismo de memoria compartida
directamente en el intérprete.

OMP4Py es una biblioteca que adapta el modelo de programación OpenMP al entorno
Python~\cite{PineiroOMP4PyFGCS2026,PineiroOMP4PyCGO2026}. Permite expresar
paralelismo mediante directivas integradas en el código, con una sintaxis próxima
a la de OpenMP en C o C++~\cite{OpenMPSpec}. La biblioteca ofrece cuatro modos
de ejecución (desde un intérprete puro hasta kernels Cython con tipos
nativos), lo que permite comparar el coste de cada capa de abstracción sobre
el mismo código fuente.

Este TFM evalúa OMP4Py sobre NAS Parallel Benchmarks adaptados a Python. Los
NAS son adecuados para este propósito porque definen problemas conocidos, clases
de tamaño estandarizadas y rutinas de verificación propias~\cite{NPB1994}. El
trabajo no consiste en validar implementaciones triviales. Se trata de portar
benchmarks completos, conservar su semántica y medir si OMP4Py puede explotarlos
de forma paralela y eficiente en un entorno de supercomputación real.


\section{Motivación}
\label{sec:motivacion}

La pregunta central del TFM es si un programador puede expresar paralelismo de
estilo OpenMP en Python con OMP4Py, obtener resultados correctos y conseguir
aceleración real sobre más de un hilo.

La motivación surge de una tensión habitual en computación científica de altas
prestaciones~\cite{IntrotoHPC2010}. Python es cómodo para desarrollar, pero los
kernels numéricos intensivos suelen requerir código nativo o bibliotecas como
NumPy~\cite{NumPy2020}, que delegan el trabajo paralelo a capas C sin exponer el
modelo de programación al desarrollador. OMP4Py propone una vía
intermedia. Mantiene una interfaz Python y, cuando el rendimiento lo exige,
compila los kernels mediante Cython. Este TFM estudia esa vía con benchmarks
estandarizados, separando el efecto del modo de ejecución del efecto del
paralelismo entre hilos.

Durante la evaluación, los dos factores no pesaron por igual. El tipado Cython
resultó decisivo, mucho más que el número de hilos, y la escalabilidad entre
hilos dependió sobre todo de la estructura de cada kernel. Las versiones de
Python \emph{free-threaded} evaluadas se comportaron de forma parecida una vez
corregida la afinidad del entorno. Las cifras concretas se presentan en los
Capítulos~\ref{chap:resultados} y~\ref{chap:versiones}.

La motivación se concreta en los objetivos de la sección siguiente y en las
contribuciones resumidas en la Sección~\ref{sec:contribuciones}.


\section{Objetivos}
\label{sec:objetivos}

El objetivo general es desarrollar y evaluar aplicaciones paralelas con OMP4Py
usando NAS Parallel Benchmarks como conjunto de referencia en el cluster
Finisterrae~III.

Los objetivos específicos se recogen en la Tabla~\ref{tab:objetivos},
identificados como O1--O7 para poder referenciarlos a lo largo de la memoria.

\begin{table}[htbp]
\centering
\small
\renewcommand{\arraystretch}{1.3}
\caption{Objetivos específicos del trabajo.}
\label{tab:objetivos}
\begin{tabular}{@{}c|p{0.85\textwidth}@{}}
\hline
\textbf{ID} & \textbf{Objetivo específico} \\
\hline
\textbf{O1} & Estudiar el modelo de programación de OMP4Py y sus cuatro modos de
      ejecución. \\
\hline
\textbf{O2} & Adaptar cinco NAS Parallel Benchmarks a OMP4Py conservando la verificación
      oficial de cada uno como criterio de aceptación. \\
\hline
\textbf{O3} & Construir una infraestructura reproducible para ejecutar los benchmarks en
      Finisterrae~III mediante Slurm. \\
\hline
\textbf{O4} & Medir rendimiento y escalabilidad con los cuatro modos de ejecución,
      distintas clases de problema y entre 1 y 32~hilos. \\
\hline
\textbf{O5} & Comparar el impacto del tipado Cython frente al aumento del número de hilos
      en cada benchmark. \\
\hline
\textbf{O6} & Comparar el comportamiento de Python~3.13t, Python~3.14t y Python~3.15t en
      modo \emph{free-threaded}. \\
\hline
\textbf{O7} & Identificar mediante profiling los cuellos de botella en los casos donde la
      escalabilidad sea reducida, usando \texttt{perf stat} y el nuevo sistema de
      profiling de Python~3.15 (muestreo y trazado) para medir el coste de
      sincronización, el recuento de operaciones del runtime y el balance de
      carga. \\
\hline
\end{tabular}
\end{table}


\section{Contribuciones}
\label{sec:contribuciones}

La memoria aporta cuatro contribuciones principales.

\begin{enumerate}
  \item Un port verificado de CG, EP, FT, IS y MG a OMP4Py, manteniendo la
        interfaz común de ejecución, las clases NAS y las rutinas internas de
        verificación.
  \item Una infraestructura reproducible de campañas en Finisterrae~III basada
        en Slurm, con generación de matrices experimentales, registros por
        ejecución y consolidación automática de resultados.
  \item Una caracterización experimental de OMP4Py con Python~3.15t en los
        cuatro modos de ejecución cuando el coste lo permite, separando el efecto
        del tipado Cython, el tamaño de clase y el número de hilos.
  \item Una comparación entre Python~3.13t, 3.14t y 3.15t, junto con un
        diagnóstico de escalabilidad mediante \texttt{perf stat} y el nuevo
        sistema de profiling de Python~3.15.
\end{enumerate}

El código de los benchmarks adaptados y de la infraestructura de campañas está
disponible en el repositorio público
\url{https://github.com/alexxvedo/omp4py_TFM}.

Estas contribuciones no se plantean como una evaluación genérica de Python
\emph{free-threaded}. El objeto de estudio es OMP4Py sobre aplicaciones NAS
completas, con especial atención a distinguir entre corrección funcional, mejora
por compilación y escalabilidad paralela efectiva.


\section{Alcance}
\label{sec:alcance}

El trabajo cubre los benchmarks EP, CG, IS, MG y FT. Son los cinco sobre los que
se ha introducido paralelismo explícito con directivas OMP4Py y cuya verificación
es correcta en todos los modos. Los benchmarks BT, SP y LU quedan fuera del
alcance principal. Sus estructuras de datos y patrones de sincronización
presentan restricciones que exceden el tiempo disponible para este TFM, y sus
adaptaciones actuales no incorporan paralelismo.

Las modificaciones a la propia biblioteca OMP4Py también quedan fuera del
alcance. La biblioteca se toma como entorno dado. El trabajo se centra en las
aplicaciones y en su evaluación experimental.

La campaña experimental principal cubre las clases de problema S, W y~A. La
clase~S permite comparar los cuatro modos de ejecución y validar la corrección
en todos los benchmarks. La clase~W se centra en los modos compilados y la
clase~A en el modo~3, que es el único modo práctico para cargas grandes. Además,
se ejecuta una campaña de comparación entre Python~3.13t, 3.14t y 3.15t sobre
las clases~S y~W. Todas las ejecuciones finales obtuvieron verificación
correcta.


\section{Estructura de la memoria}
\label{sec:estructura}

El Capítulo~\ref{chap:contexto} describe el contexto tecnológico, incluyendo Python
\emph{free-threaded}, OMP4Py y los NAS Parallel Benchmarks. El
Capítulo~\ref{chap:adaptacion-nas} detalla el proceso de adaptación de los cinco
benchmarks y las decisiones de implementación más relevantes. El
Capítulo~\ref{chap:metodologia} define la metodología experimental y la
infraestructura de medición. El Capítulo~\ref{chap:resultados} presenta y analiza
los resultados de la campaña principal con Python~3.15t. El
Capítulo~\ref{chap:profiling} diagnostica con \texttt{perf stat} la ocupación
real de CPU, el IPC y los límites de escalabilidad, y con el nuevo sistema de
profiling de Python~3.15 (muestreo y trazado) los costes de sincronización, el
recuento de operaciones del runtime y el balance de carga. El
Capítulo~\ref{chap:versiones} compara Python~3.13t, 3.14t y 3.15t. El
Capítulo~\ref{chap:conclusiones} resume las principales conclusiones y plantea el trabajo
futuro.
